Table~\ref{tab:llm-controllers} shows the conventional alternative to a small router: a small instruction-tuned LLM
driving the MCP tools end to end. It is served locally and issues JSON actions through the repository's
\texttt{LocalModel} adapter.

Qwen2.5-0.5B-Instruct (494M parameters, measured by the server) accepted none of the six search-arm episodes it ran
on MCP-02 test worlds. It used nearly the whole 40-call budget and processed about 79K prompt and completion tokens
per episode, yet never submitted a packet. In five episodes it called \texttt{resource\_read} 40 times with invented
URIs, never searching, and every call returned an error. In the sixth it repeated a nonexistent tool name.

The four-pass looped router solves all 32 episodes of the same cell in 5.5 calls. It scores about 2.7K encoder
positions per episode, in 0.14 s of GPU time.

The comparison is lopsided by design: the router works inside the shared plan tracker and frontier, while the LLM
must also plan, parse and compose. That is the point. Spending local compute on a general small LLM is the
expensive way to explore MCP surfaces.

We stopped the LLM stage after six episodes to stay within the VRAM lease on a GPU shared with other workloads. The
planned Qwen3-1.7B run was not executed.

\IfFileExists{tables/llm_controllers.tex}{\begin{table}[t]
\centering
\small
\begin{adjustbox}{max width=\linewidth}
\begin{tabular}{lrrrrr}
\toprule
Controller & Arm & Accepted & Calls/ep & Local tokens/ep & Local s/ep \\
\midrule
llm:qwen2.5-0.5b-instruct & search & 0/6 & 39.83 & 78520.33 & 334.74 \\
lexical\_bm25 & search & 32/32 & 7.22 & 0.00 & 0.01 \\
looped\_p4\_d128 & search & 32/32 & 5.53 & 2680.47 & 0.14 \\
trm\_c2i1\_d128 & search & 32/32 & 6.25 & 3676.25 & 0.49 \\
\bottomrule
\end{tabular}
\end{adjustbox}
\caption{Local LLMs as end-to-end MCP controllers (JSON actions through MCP.local\_model.LocalModel, RTX 5080) against small routers on MCP-02 test worlds at 256 resources. LLM tokens are server-counted prompt plus completion tokens; router tokens are encoder positions scored. Router rows cover 32 worlds, LLM rows 8.}
\label{tab:llm-controllers}
\end{table}
}{}
